\documentclass[12pt]{article}
\usepackage[utf8]{inputenc}
\usepackage[T1]{fontenc}
\usepackage{hyperref}
\title{Zaniza Zapotec Language Resources for K-12 Curriculum Development}
\author{Compiled by Claude Sonnet 5.5 (claude-sonnet-5-5) for Andr\'as Kornai}
\date{}
\begin{document}
\maketitle

\section{Introduction}
Zaniza Zapotec (ISO 639-3 zpw; Glottocode zani1235) is an endangered Otomanguean language of Santa Mar\'ia Zaniza and Santiago Textitl\'an, Oaxaca, Mexico. It is also called Western Sola de Vega Zapotec and belongs to the Papabuco group. Wikipedia gives about 400 fluent speakers (2003), while the survey brief says about 770. Very little material exists: one linguistic description and evangelistic audio.

\section{Primary Educational Resources}
No primers, textbooks or school materials were found.

\section{Linguistic and Reference Materials}
\begin{itemize}
\item Wikipedia article on Zaniza Zapotec, including phonology and references: \url{https://en.wikipedia.org/wiki/Zaniza_Zapotec}
\end{itemize}
The Wikipedia references cite Natalie Operstein, a 2002 paper on positional verbs and relational nouns (Santa Barbara Papers in Linguistics, vol. 11), and a 2015 book on Zaniza Zapotec. I did not see direct links to either.

\section{Religious and Community Resources}
No scripture translation was found.

\section{Digital and Audio Resources}
\begin{itemize}
\item Global Recordings Network, audio (MP3 and slideshow packages): \url{https://globalrecordings.net/en/language/zpw}
\end{itemize}

\section{Curriculum Development Context}
The language is used by all adults but not all young people, and no school use is known. There is no dictionary or reading material in the sources found, so a curriculum would need basic lexical and orthographic work with the community. Operstein's grammatical description is the main academic starting point.

\end{document}